% Chinese working edition. Original book class, geometry, headings and Latin font retained.
\documentclass[a4paper,oneside,fleqn]{book}
\usepackage{geometry}
% Upstream running heads use two lines; preserve the text block's top position.
\setlength{\headheight}{22pt}
\addtolength{\topmargin}{-10pt}
\PassOptionsToPackage{no-math}{fontspec}
\usepackage[fontset=none,scheme=plain]{ctex}
\setCJKmainfont{Noto Serif CJK SC}[ItalicFont=Noto Serif CJK SC,ItalicFeatures={FakeSlant=0.2}]
\setCJKsansfont{Noto Sans CJK SC}
\setCJKmonofont{Noto Sans Mono CJK SC}
\def\mode{WebMode}
\usepackage{ifthen}
%!TEX root = main.tex
\usepackage[T1]{fontenc}
\usepackage[final]{microtype}
\usepackage{amsmath,amssymb}
\usepackage{amsthm,mathtools}
\usepackage{amscd}
\usepackage{etoolbox}
\usepackage[cal=esstix,scr=rsfs]{mathalfa}



% DOUBLE QUOTES
\usepackage[autostyle=false,style=english]{csquotes}
\MakeOuterQuote{"}



% HEADERS AND FOOTERS
\usepackage{fancyhdr}
\renewcommand{\chaptermark}[1]{\markboth{\thechapter. #1}{}}
\renewcommand{\sectionmark}[1]{\markboth{\thesection. #1}{}}
\fancyhf{}
\ifthenelse{\equal{\mode}{WebMode}}{%
  \fancyhead[C]{\small FGA~\thepart\\\leftmark}%
  \renewcommand{\headrulewidth}{0.5pt}%
}{%
  \fancyhead[LO]{FGA~\thepart}%
  \fancyhead[RE]{\leftmark}%
  \renewcommand{\headrulewidth}{0pt}%
}
\fancyfoot[C]{\thepage}



% BIBLIOGRAPHY
\usepackage[backend=bibtex]{biblatex}
\renewbibmacro{in:}{%
  \ifboolexpr{%
     test {\ifentrytype{article}}%
  }{}{\printtext{\bibstring{in}\intitlepunct}}%
}



% LINKS
\ifthenelse{\equal{\mode}{WebMode}}{%
  \usepackage{xcolor}%
  \usepackage[colorlinks=true,linkcolor={blue!60!black},citecolor={blue!60!black},urlcolor={purple!60!black}]{hyperref}%
}{%
  \usepackage[colorlinks=false,hidelinks]{hyperref}
}

\usepackage[nameinlink,capitalize]{cleveref}
\crefname{section}{§}{§§}
\crefname{subsection}{§}{§§}
\labelformat{subsection}{\thesection.#1}



% HEADINGS
\usepackage{titlesec}

\titleformat{\part}[display]{}{\huge FGA~\thepart}{8pt}{\huge}
\titleclass{\part}{top}
\titlespacing{\part}{0pt}{0pt}{*6}

\titleformat{\chapter}[hang]{}{\Large\thechapter.}{8pt}{\Large}
\ifthenelse{\equal{\mode}{WebMode}}{%
  \titleclass{\chapter}{straight}%
  \titlespacing{\chapter}{0pt}{*8}{*4}%
}{%
  \titlespacing{\chapter}{0pt}{0pt}{*4}%
}

\titleformat{\section}[hang]{}{\large\thesection.}{8pt}{\large}
\titlespacing{\section}{0pt}{*8}{*2}


\titleformat{\subsection}[hang]{}{\large\thesection.\thesubsection.}{8pt}{\large}
\titlespacing{\subsection}{0pt}{*8}{*2}




% ENVIRONMENTS
\makeatletter
\def\skipsinglepar{%
  \@ifnextchar\par
    {\expandafter\skipsinglepar\@gobble}%
    {}%
}
\makeatother


%% Abstract
\newenvironment{abstract}{\par\vspace{\topsep}\noindent\textbf{译者说明。}~}{\par\vspace{\topsep}}


%% Numbering for \ref
\makeatletter
\newcommand{\setnumber}[1]{\def\@currentlabel{#1}}
\makeatother


%% Numbered equations
\newenvironment{eqn}[1]{\vspace{-1em}}{}


%% Original citations
\newenvironment{original}{\par\vspace{\topsep}\hrule\vspace{\topsep}\noindent}{\vspace{\topsep}\hrule\par\vspace{\topsep}}


%% Custom inner environment
\newenvironment{inenv}[4][]{%
  % first argument (optional) is the title
  % second argument (mandatory) is the number
  % third argument (mandatory) is the taxon
  % fourth argument (mandatory) is empty for non-italic
  \par\vspace{\topsep}\noindent%
  \ifstrempty{#2}{%if no number
    \ifstrempty{#1}{%if no title
      \textbf{#3.}%
    }{%if title
      \textbf{#3} (#1).%
    }%
  }{%if number
    \ifstrempty{#1}{%
      \textbf{#3~#2.}%
    }{%
      \textbf{#3~#2} (#1).%
    }%
  }%
  \ifstrempty{#4}{}{\itshape}~}%non-italic
  {\par\vspace{\topsep}}


%% Italicised environments
\newenvironment{theorem}[2][]%
  {\begin{inenv}[#1]{#2}{定理}{italic}}%
  {\end{inenv}}
\newenvironment{lemma}[2][]%
  {\begin{inenv}[#1]{#2}{引理}{italic}}%
  {\end{inenv}}
\newenvironment{proposition}[2][]%
  {\begin{inenv}[#1]{#2}{命题}{italic}}%
  {\end{inenv}}
\newenvironment{corollary}[2][]%
  {\begin{inenv}[#1]{#2}{推论}{italic}}%
  {\end{inenv}}
\newenvironment{statement}[2][]%
  {\begin{inenv}[#1]{#2}{断言}{italic}}%
  {\end{inenv}}
\newenvironment{conjecture}[2][]%
  {\begin{inenv}[#1]{#2}{猜想}{italic}}%
  {\end{inenv}}


%% Upright environments
\newenvironment{definition}[2][]%
  {\begin{inenv}[#1]{#2}{定义}{}}%
  {\end{inenv}}
\newenvironment{remark}[2][]%
  {\begin{inenv}[#1]{#2}{注}{}}%
  {\end{inenv}}
\newenvironment{remarks}[2][]%
  {\begin{inenv}[#1]{#2}{注}{}}%
  {\end{inenv}}
\newenvironment{example}[2][]%
  {\begin{inenv}[#1]{#2}{例}{}}%
  {\end{inenv}}
\newenvironment{scholium}[2][]%
  {\begin{inenv}[#1]{#2}{评注}{}}%
  {\end{inenv}}

% Keep the upstream book's layout and labels; localize only reader-facing names.
\renewcommand{\proofname}{证明}
\AtEveryBibitem{\label{\thefield{entrykey}}}

\renewcommand{\contentsname}{目录}

\usepackage{Baskervaldx}
\AtEndPreamble{\setmainfont{Baskervaldx-Reg.otf}[BoldFont=Baskervaldx-Bol.otf,ItalicFont=Baskervaldx-Ita.otf,BoldItalicFont=Baskervaldx-BolIta.otf]}
\special{pdf:mapfile Baskervaldx.map}
\special{pdf:mapfile ESSTIX.map}
\addbibresource{fga-bibliography.bib}
\title{代数几何基础\\[0.5em]\large Fondements de la Géométrie Algébrique\\[1em]\normalsize 中文工作稿：前言与报告 I 第 1--2 节}
\author{Alexander Grothendieck}
\date{1956--1962}
\begin{document}
\frontmatter
\pagestyle{plain}
\maketitle
\begin{abstract}
本稿据 Tim Hosgood 的英文原生稿翻译，并在发现疑误处核对法文原本。本次仅含前言、第一篇报告的引言及第 1--2 节；其他七篇报告和第一篇其余各节尚未译入。原书作者、报告编号、公式与引用均予保留。英文稿与法文之间的差异及校核状态另见源码包中的 SOURCE\_NOTES.json。本稿尚未完成独立数学审校。
\end{abstract}
\setcounter{tocdepth}{1}
\tableofcontents
\mainmatter
\pagestyle{fancy}
\setcounter{part}{-1}
\renewcommand\thepart{}
\part{前言}\label{fga-foreword}\thispagestyle{fancy}\par{}1957 至 1962 年间，我们在 \emph{Séminaire Bourbaki} 上作了八次关于代数几何基础的报告。除第一篇外，这些报告均使用\emph{概形}的语言。报告中陈述的所有结果都将在 Jean Dieudonné 与 Alexander Grothendieck 的《Éléments de Géométrie Algébrique》中得到论述。不过，这些报告的主要内容，目前既不见于 EGA 已经出版或仍在准备中的任何一章，也不见于其他任何书籍或论文；这种情形很可能还会持续几年。这是促使我们将这些报告汇集成册的主要原因：在一份完善的综述问世之前，先让读者能够接触概形理论中的若干思想和关键结果。阅读这些报告，还能使读者迅速熟悉上述结果和思想，而不必先处理系统论述中难以避免的繁复细节；同时，也能为进一步研读这种系统论述提供必要的动机。\par{}为方便读者，我们在这里按节汇集了一些评注和勘误，特别说明各篇写成以来取得的进展，并补充若干参考文献。\emph{[Trans.] 我们没有在此集中翻译这些评注和勘误，而是将它们插入正文的相应位置；除已直接改正而未另作标记的细小订正外，均以“\emph{[Comp.]}”开头，表示“Complément”。}\par{}这些文章中的许多结果，已在 \emph{Séminaire de Géométrie Algébrique du Bois Marie} 以及后来于 1961 至 1962 年在 Harvard 举行的两个讨论班中得到详细讨论。后两个讨论班的第一个由我主持，第二个由 Mumford–Tate 主持，其笔记目前正由 Lichtenbaum 整理。\emph{[Trans.] 这些报告后来还多次提及 Mumford–Tate 的笔记，但这份笔记实际上并不存在：它从未完成，也没有任何草稿发表过（参见 \href{https://mathoverflow.net/a/452230/73622}{MO:452230/73622}）。}

\clearpage
\renewcommand\thepart{\arabic{part}}
\part{凝聚代数层的对偶定理}\label{fga1}\thispagestyle{fancy}\begin{original}{}\label{fga1-original-citation}\skipsinglepar\par{}A. Grothendieck.
“Théorème de dualité pour les faisceaux algébriques cohérents”.
\emph{Séminaire Bourbaki} \textbf{9}（1956–57），第 149 号报告。
\url{http://www.numdam.org/book-part/SB_1956-1958__4__169_0/}\end{original}\par{}下面的结果受到 Serre 的“代数对偶定理”的启发，发现于 1955 年冬季和 1956 年冬季。利用射影空间上同调的一些相当初等的结果 \hyperref[Ser1955]{[Ser1955]}，并充分使用 Cartan–Eilenberg 同调代数在 \hyperref[Gro1957]{[Gro1957]} 中给出的形式，就能很简便地建立这些结果。\chapter{模层的 \texorpdfstring{\(\operatorname {Ext}\)}{Ext}}\thispagestyle{fancy}\label{fga1-1}\par{}\emph{（参见 [\hyperref[Gro1957]{Gro1957}，第 3、4 章]）}\par{}设 \(X\) 是一个拓扑空间，其上给定一个有单位元的环层 \(\mathcal {O}\)，环不必交换。考虑由 \(\mathcal {O}\)-模层组成的 Abel 范畴 \(\mathcal {C}^\mathcal {O}\)，这些模层也简称为 \(\mathcal {O}\)-模。已知这个范畴中的每个对象都有内射分解，因此可以定义具有熟知形式性质的 \(\operatorname {Ext}\) 函子。为避免混淆，更确切地说，我们用 \(\operatorname {Hom}_\mathcal {O}(X;\mathscr {A},\mathscr {B})\)，或简记为 \(\operatorname {Hom}(X;\mathscr {A},\mathscr {B})\)，表示从 \(\mathscr {A}\) 到 \(\mathscr {B}\) 的 \(\mathcal {O}\)-同态所成的 Abel \emph{群}；而 \(\mathscr {H}\kern  -2.5pt\mathit {om}_\mathcal {O}(\mathscr {A},\mathscr {B})\) 表示从 \(\mathscr {A}\) 到 \(\mathscr {B}\) 的同态芽所成的\emph{层}（其中 \(\mathscr {A},\mathscr {B}\in \mathcal {C}^\mathcal {O}\)）。固定 \(\mathscr {A}\in \mathcal {C}^\mathcal {O}\) 后，定义函子 \(h_\mathscr {A}\) 和 \(\underline {h}_\mathscr {A}\)，它们分别取值于 Abel 群范畴 \(\mathcal {C}\) 和 \(X\) 上的 Abel 层范畴 \(\mathcal {C}^Z\)，定义式为：
\begin{eqn}{1.1}\setnumber{1.1}\label{fga1-equation-1.1}\skipsinglepar\[       \begin {aligned}         h_\mathscr {A}(\mathscr {B})         &= \operatorname {Hom}_\mathcal {O}(X;\mathscr {A},\mathscr {B})       \\\underline {h}_\mathscr {A}(\mathscr {B})         &= \mathscr {H}\kern  -2.5pt\mathit {om}_\mathcal {O}(\mathscr {A},\mathscr {B}).       \end {aligned}     \tag{1.1}     \]\end{eqn}
函子 \(h_\mathscr {A}\) 和 \(\underline {h}_\mathscr {A}\) 都是左正合协变函子，因此可以考虑其右导出函子，分别记为 \(\operatorname {Ext}_\mathcal {O}^p(X;\mathscr {A},\mathscr {B})\) 和 \(\mathscr {E}\kern  -1.5pt\mathit {xt}_\mathcal {O}^p(\mathscr {A},\mathscr {B})\)。按照定义，有
\begin{eqn}{1.2}\setnumber{1.2}\label{fga1-equation-1.2}\skipsinglepar\[       \begin {gathered}         \operatorname {Ext}_\mathcal {O}^p(X;\mathscr {A},\mathscr {B})         = (\operatorname {R}^p h_\mathscr {A})(\mathscr {B})         = \operatorname {H}^p(\operatorname {Hom}_\mathcal {O}(X;\mathscr {A},C(\mathscr {B})))       \\\mathscr {E}\kern  -1.5pt\mathit {xt}_\mathcal {O}^p(\mathscr {A},\mathscr {B})         = (\operatorname {R}^p \underline {h}_\mathscr {A})(\mathscr {B})         = \operatorname {H}^p(\mathscr {H}\kern  -2.5pt\mathit {om}_\mathcal {O}(\mathscr {A},C(\mathscr {B})))       \end {gathered}     \tag{1.2}     \]\end{eqn}
其中 \(\operatorname {R}^p\) 表示取右导出函子，\(C(\mathscr {B})\) 表示 \(\mathscr {B}\) 在 \(\mathcal {C}^\mathcal {O}\) 中的任意内射分解。用 \(\Gamma \colon \mathcal {C}^Z\to \mathcal {C}\) 表示“截面”函子；回顾其右导出函子记为 \(B\mapsto \operatorname {H}^p(X,\mathscr {B})\)：
\begin{eqn}{1.3}\setnumber{1.3}\label{fga1-equation-1.3}\skipsinglepar\[       \operatorname {H}^p(X,\mathscr {B})       = (\operatorname {R}^p\Gamma )(\mathscr {B})       = \operatorname {H}^p(\Gamma (C(\mathscr {B}))).     \tag{1.3}     \]\end{eqn}
显然有 \(h_\mathscr {A}=\Gamma \underline {h}_\mathscr {A}\)；此外还可以证明，\(\underline {h}_\mathscr {A}\) 将内射对象映为 \(\Gamma \)-无环对象。由此可按熟知的方法得到：\begin{proposition}{1}\setnumber{1}\label{fga1-proposition-1}\skipsinglepar\par{}对每个 \(\mathcal {O}\)-模 \(\mathscr {A}\)，在 \(\mathcal {C}^\mathcal {O}\) 上都存在一个上同调谱函子，收敛于分次函子 \((\operatorname {Ext}_\mathcal {O}^\bullet (X;\mathscr {A},\mathscr {B}))\)，其初始页为
\begin{eqn}{1.4}\setnumber{1.4}\label{fga1-equation-1.4}\skipsinglepar\[         E_2^{p,q}(\mathscr {A},\mathscr {B})         = \operatorname {H}^p(X,\mathscr {E}\kern  -1.5pt\mathit {xt}_\mathcal {O}^q(\mathscr {A},\mathscr {B})).       \tag{1.4}       \]\end{eqn}\end{proposition}\par{}由此得到“\emph{边缘同态}”，以及一个我们不再写出的五项正合列。\footnote{中译校记：法文原本第 149-02 页为五项正合列；英文稿此处误作短正合列。}\begin{corollary}{1}\setnumber{1}\label{fga1-proposition-1-corollary-1}\skipsinglepar\par{}如果 \(\mathscr {A}\) 局部同构于 \(\mathcal {O}^n\)，则有典范同构
\begin{eqn}{1.5}\setnumber{1.5}\label{fga1-equation-1.5}\skipsinglepar\[         \operatorname {Ext}_\mathcal {O}^p(X;\mathscr {A},\mathscr {B})         \xleftarrow {\sim } \operatorname {H}^p(X,\mathscr {H}\kern  -2.5pt\mathit {om}_\mathcal {O}(\mathscr {A},\mathscr {B}))       \tag{1.5}       \]\end{eqn}
（由谱序列的边缘同态给出）。\footnote{中译校记：式 (1.5) 中的空间依照法文原本记为 $X$；英文稿误为小写 $x$。}特别地，有典范同构
\begin{eqn}{1.6}\setnumber{1.6}\label{fga1-equation-1.6}\skipsinglepar\[         \operatorname {Ext}_\mathcal {O}^p(X;\mathcal {O},\mathscr {B})         = \operatorname {H}^p(X,\mathscr {B}).       \tag{1.6}       \]\end{eqn}\end{corollary}\par{}为使用这些结果，需要知道如何具体描述 \(\operatorname {Ext}_\mathcal {O}^p(\mathscr {A},\mathscr {B})\)。这些函子可以局部计算，也就是说，如果 \(U\) 是 \(X\) 的一个开子集，则
\[     \mathscr {E}\kern  -1.5pt\mathit {xt}_\mathcal {O}^p(\mathscr {A},\mathscr {B})|U     = \mathscr {E}\kern  -1.5pt\mathit {xt}_{\mathcal {O}|U}^p(\mathscr {A}|U,\mathscr {B}|U)   \]
这是因为内射 \(\mathcal {O}\)-模在 \(U\) 上的限制仍是内射 \((\mathcal {O}|U)\)-模。此外，固定 \(x\in  X\) 后，有函子性的同态
\begin{eqn}{1.7}\setnumber{1.7}\label{fga1-equation-1.7}\skipsinglepar\[       \mathscr {H}\kern  -2.5pt\mathit {om}_\mathcal {O}(\mathscr {A},\mathscr {B})_x       \to  \operatorname {Hom}_{\mathcal {O}_x}(\mathscr {A}_x,\mathscr {B}_x)     \tag{1.7}     \]\end{eqn}
它们唯一地延拓为关于 \(\mathscr {B}\) 的上同调函子之间的同态：
\begin{eqn}{1.8}\setnumber{1.8}\label{fga1-equation-1.8}\skipsinglepar\[       \mathscr {E}\kern  -1.5pt\mathit {xt}_\mathcal {O}^p(\mathscr {A},\mathscr {B})_x       \to  \operatorname {Ext}_{\mathcal {O}_x}^p(\mathscr {A}_x,\mathscr {B}_x).     \tag{1.8}     \]\end{eqn}\begin{proposition}{2}\setnumber{2}\label{fga1-proposition-2}\skipsinglepar\par{}如果 \(\mathscr {A}\) 在 \(x\) 的某个邻域内同构于某个同态 \(\mathcal {O}^m\to \mathcal {O}^n\) 的余核，那么对于所有 \(p\)，\hyperref[fga1-equation-1.7]{式~\ref*{fga1-equation-1.7}} 都是同构。特别地，当 \(\mathscr {A}\) 是凝聚 \(\mathcal {O}\)-模 \hyperref[Ser1955]{[Ser1955]} 时，情形就是如此。\end{proposition}\begin{proposition}{3}\setnumber{3}\label{fga1-proposition-3}\skipsinglepar\par{}设 \(\mathscr {L}_\bullet =(\mathscr {L}_i)\) 是 \(\mathcal {O}\)-模 \(\mathscr {A}\) 的一个左分解，其各项都是局部同构于某个 \(\mathcal {O}^n\) 的 \(\mathcal {O}\)-模。那么 \(\mathscr {E}\kern  -1.5pt\mathit {xt}_\mathcal {O}(\mathscr {A},\mathscr {B})\) 可与 \(\operatorname {H}^\bullet (\mathscr {H}\kern  -2.5pt\mathit {om}_\mathcal {O}(\mathscr {L}_\bullet ,\mathscr {B}))\) 等同，而 \(\operatorname {Ext}_\mathcal {O}(X;\mathscr {A},\mathscr {B})\) 可与 \(X\) 关于复形 \(\mathscr {H}\kern  -2.5pt\mathit {om}_\mathcal {O}(\mathscr {L}_\bullet ,\mathscr {B})\) 的超上同调等同。\begin{proof}\par{}这是标准的论证：考虑双复形 \(\mathscr {H}\kern  -2.5pt\mathit {om}_\mathcal {O}(\mathscr {L}_\bullet ,C(\mathscr {B}))\)，其中 \(C(\mathscr {B})\) 是 \(\mathscr {B}\) 的内射分解，并考虑从 \(\mathscr {H}\kern  -2.5pt\mathit {om}_\mathcal {O}(\mathscr {L}_\bullet ,\mathscr {B})\) 和 \(\mathscr {H}\kern  -2.5pt\mathit {om}_\mathcal {O}(\mathscr {A},C(\mathscr {B}))\) 到这个双复形的自然同态。\end{proof}\end{proposition}\par{}最后指出，\hyperref[fga1-equation-1.2]{式~\ref*{fga1-equation-1.2}} 中引入的两个 \(\operatorname {Ext}\) 函子，不仅是关于 \(\mathscr {B}\) 的上同调函子，而且是\emph{上同调双函子}，关于 \(\mathscr {B}\) 协变，关于 \(\mathscr {A}\) 反变。\chapter{\texorpdfstring{\(\operatorname {Ext}\)}{Ext} 中的复合律}\thispagestyle{fancy}\label{fga1-2}\par{}本节的结果由 Cartier 和 Yoneda 各自独立得到；详情可参见 Cartier 的报告 \hyperref[Car1957]{[Car1957]}。设 \(\mathcal {C}\) 是 Abel 范畴，\(K\) 和 \(L\) 是 \(\mathcal {C}\) 中的两个分次对象。用 \(\operatorname {Hom}(K,L)\) 表示这样一个分次 Abel 群：其 \(n\) 次分量由从 \(K\) 到 \(L\) 的 \(n\) 次齐次同态组成，也就是由同态 \(K^i\to  L^{i+n}\) 组成的族 \((u_i)\)。如果 \(K\) 和 \(L\) 是复形（为固定约定，设微分的次数为 \(+1\)），则在 \(\operatorname {Hom}(K,L)\) 上引入如下微分算子：
\begin{eqn}{2.1}\setnumber{2.1}\label{fga1-equation-2.1}\skipsinglepar\[       \delta (u) = \mathrm {d}u + (-1)^{n+1}u\mathrm {d}       \quad \text {where }n=\deg (u)     \tag{2.1}     \]\end{eqn}
这使它成为一个微分次数为 \(+1\) 的复形。\(n\) 次闭链就是作为齐次映射与 \(u\) 反交换的 \(n\) 次映射。于是可以考虑 \(\operatorname {H}^\bullet (\operatorname {Hom}(K,L))\)；它是 \(K\) 与 \(L\) 的同伦型的不变量，有时记为 \(\operatorname {H}^\bullet (K,L)\)。如果还有第三个复形 \(M\)，则同态的复合定义一个与微分算子相容的配对 \(\operatorname {Hom}(K,L)\times \operatorname {Hom}(L,M)\to \operatorname {Hom}(K,M)\)，再取上同调，就得到配对
\begin{eqn}{2.2}\setnumber{2.2}\label{fga1-equation-2.2}\skipsinglepar\[       \operatorname {H}^\bullet (K,L)\times \operatorname {H}^\bullet (L,M) \to  \operatorname {H}^\bullet (K,M)     \tag{2.2}     \]\end{eqn}
记作 \((u,v)\mapsto  vu\)。这些配对显然满足结合律；特别地，\(\operatorname {H}^\bullet (K,K)\) 是一个有单位元的结合分次环，而 \(\operatorname {H}^\bullet (K,L)\)（分别地，\(\operatorname {H}^\bullet (L,K)\)）是这个环上的分次右模（分别地，左模），等等。在 \(0\) 次时，\hyperref[fga1-equation-2.2]{式~\ref*{fga1-equation-2.2}} 就是复形的容许同态的复合。最后，若复形正合列 \(0\to  K'\to  K\to  K''\to 0\) 满足：对每个 \(i\)，\(K'^i\) 可与 \(K^i\) 的一个直和因子等同，则得到群复形 \(\operatorname {Hom}(K'',L)\) 等所成的正合列，从而得到连接算子 \(\operatorname {H}^i(K',L)\to \operatorname {H}^{i+1}(K'',L)\)。同样，可以定义与 \(L\) 中的一个正合列相联系的连接算子。\hyperref[fga1-equation-2.2]{式~\ref*{fga1-equation-2.2}} 中的配对按通常的意义与这些连接算子相容。\par{}现在设 \(\mathcal {C}\) 是这样一个范畴：\(\mathcal {C}\) 中的每个对象 \(A\) 都有一个内射分解 \(C(A)\)。利用双复形定理的许多变体之一，可以看到
\[     \operatorname {H}^\bullet (C(A),C(B)) = \operatorname {H}^\bullet (\operatorname {Hom}(C(A),C(B)))   \]
典范同构于
\[     \operatorname {H}^\bullet (\operatorname {Hom}(A,C(B))) = \operatorname {Ext}^\bullet (A,B).   \]
上面讨论的连接算子给出 \(\operatorname {Ext}\) 的连接同态。此外，\hyperref[fga1-equation-2.2]{式~\ref*{fga1-equation-2.2}} 中的配对在这里给出结合的配对
\begin{eqn}{2.3}\setnumber{2.3}\label{fga1-equation-2.3}\skipsinglepar\[       \operatorname {Ext}^\bullet (A,B)\times \operatorname {Ext}^\bullet (B,C) \to  \operatorname {Ext}^\bullet (A,C)     \tag{2.3}     \]\end{eqn}
并且与连接同态相容。特别地，\(\operatorname {Ext}^\bullet (A,A)\) 是有单位元的结合分次环，等等。（同样可以证明，\(\operatorname {Ext}\) 函子作用于任意函子的导出函子；这里不使用这一事实。）\par{}如果所考虑的范畴是 \(X\) 上的 \(\mathcal {O}\)-模范畴 \(\mathcal {C}^\mathcal {O}\)，便得到配对
\begin{eqn}{2.4}\setnumber{2.4}\label{fga1-equation-2.4}\skipsinglepar\[       \operatorname {Ext}_\mathcal {O}^p(X;\mathscr {A},\mathscr {B})\times \operatorname {Ext}_\mathcal {O}^q(X;\mathscr {B},\mathscr {C}) \to  \operatorname {Ext}_\mathcal {O}^{p+q}(X;\mathscr {A},\mathscr {C})     \tag{2.4}     \]\end{eqn}
并可按照已经说明的方法将它们具体写出。用同样的方法，将 Abel 群范畴换成 \(X\) 上的 Abel 层范畴，并将 \(\operatorname {Hom}\) 函子换成 \(\mathscr {H}\kern  -2.5pt\mathit {om}\) 函子，就又能定义具有相同形式性质的配对，不过这次具有“局部性质”：
\begin{eqn}{2.5}\setnumber{2.5}\label{fga1-equation-2.5}\skipsinglepar\[       \mathscr {E}\kern  -1.5pt\mathit {xt}_\mathcal {O}^p(\mathscr {A},\mathscr {B})\times \mathscr {E}\kern  -1.5pt\mathit {xt}_\mathcal {O}^q(\mathscr {B},\mathscr {C}) \to  \mathscr {E}\kern  -1.5pt\mathit {xt}_\mathcal {O}^{p+q}(\mathscr {A},\mathscr {C}).     \tag{2.5}     \]\end{eqn}
注意到 \hyperref[fga1-equation-1.8]{式~\ref*{fga1-equation-1.8}} 中的同态与 \(\operatorname {Ext}\) 之间的配对相容，就能进一步说明这些配对。\par{}最后，回顾函子 \(\operatorname {H}^p(X,A)\) 之间还存在一种乘法结构，即杯积。于是可以看到，\hyperref[fga1-proposition-1]{命题~\ref*{fga1-proposition-1}} 的谱序列与这些乘法结构相容；更确切地说，谱序列 \(E(A,B)\) 与谱序列 \(E(B,C)\) 之间存在一个取值于谱序列 \(E(A,C)\) 的配对，它收敛于整体 \(\operatorname {Ext}\) 之间的配对，其初始页来自杯积及 \hyperref[fga1-equation-1.4]{式~\ref*{fga1-equation-1.4}} 右端局部 \(\operatorname {Ext}\) 之间的配对。特别地，由此得到“边缘同态”
\begin{eqn}{2.6}\setnumber{2.6}\label{fga1-equation-2.6}\skipsinglepar\[       \operatorname {Ext}_\mathcal {O}^n(X;\mathscr {A},\mathscr {B}) \to  \operatorname {H}^0(X;\mathscr {E}\kern  -1.5pt\mathit {xt}_\mathcal {O}^n(\mathscr {A},\mathscr {B}))     \tag{2.6}     \]\end{eqn}
\begin{eqn}{2.7}\setnumber{2.7}\label{fga1-equation-2.7}\skipsinglepar\[       \operatorname {H}^n(X,\mathscr {H}\kern  -2.5pt\mathit {om}_\mathcal {O}(\mathscr {A},\mathscr {B})) \to  \operatorname {Ext}_\mathcal {O}^n(X;\mathscr {A},\mathscr {B})     \tag{2.7}     \]\end{eqn}
与乘法结构相容。因此，如果只考虑局部同构于某个 \(\mathcal {O}^m\) 的层，结合 \hyperref[fga1-equation-1.5]{式~\ref*{fga1-equation-1.5}} 中的同构，便能用杯积完全说明整体 \(\operatorname {Ext}\) 的复合。

\backmatter
\pagestyle{plain}
\nocite{Car1957,Gro1957,Ser1955}
\printbibliography[heading=bibintoc,title=参考文献]
\end{document}
